% Section 9 of the manuscript: runtime cost (question 5). The method is stated first, as the outline requires; it
% follows sections 6.10 and 10.9 of the frozen registration and the description of the samplers in kit/e3/README.md.
% The results are written from the result files e3_macos.txt and e3_windows.txt after the analysis.

\section{Results: runtime cost}
\label{sec:cost}

\subsection{Method}
\label{sec:cost-method}

The cost of the logger to the machine was measured on each system under two conditions, one night each, for eight
hours at an interval of 30~s: \emph{idle}, in which the logger ran and neither a person nor a plan operated the machine
(what else ran on it is given in Section~\ref{sec:results-attempts}), and
\emph{working}, in which the logger ran while the runner played the plan of one agreement cell (on macOS the plan with
live conversion on, which takes about 11.9~hours; on Windows the plan of the Windows cell, which takes 12 to
13~hours). In the working condition only the samples of the first eight hours were used. The conditions were measured
after the agreement cells of the same machine (Section~\ref{sec:method-cells}).

A sampler found the logger's process by its name and recorded, at every interval, its CPU and its resident memory; no
other process was listed and no window title was read. CPU is expressed on the basis that one core is 100\%. The two
systems measure differently. On macOS the CPU value of a sample is the increase of the process's cumulative CPU time
divided by the increase of the wall-clock time, a mean over the whole interval, and the memory is the resident memory
that the system's process monitor reports. On Windows the CPU value is the processor time of the process over about
one second at the moment of reading, so a burst between two samples is not seen, and the memory is the working set.
For this reason the values are described within each system and are not compared between the systems.

Reported for each condition are the mean, median, 95th percentile and maximum of the CPU, the final resident memory,
and the least-squares slope of the resident memory in megabytes per hour. The two conditions are always reported
together and are not reduced to one figure. A condition with fewer than 480 usable samples or a window shorter than
four hours was insufficient and was run again on another night, as was a condition that was aborted, up to three
attempts in all; every attempt is reported. In the working condition the runs were checked by the stopping rule of the
agreement cells (Section~\ref{sec:method-checks}). On Windows a counter that cannot be read gives zero, so a working
condition in which every CPU sample was zero was treated as insufficient, because the logger processes the presses of
the plan; in the idle condition a true zero is possible, and such a condition is reported with this flag. Battery use
was not measured. The samples of both conditions are deposited; the observation log of the idle condition is not,
because the logger was started before the command that begins the sampling was typed, so that log holds a person's
key presses.

\subsection{Results}
\label{sec:cost-results}

The records of the runs hold one attempt of each condition on each system (the working conditions are in
Table~\ref{tab:attempts}), and every one had enough samples to be reported (Table~\ref{tab:cost}). On macOS the mean CPU of the
logger was \rv{macos.e3.idle.cpu.mean} in the idle condition and \rv{macos.e3.working.cpu.mean} in the working
condition; on Windows it was \rv{windows.e3.idle.cpu.mean} idle and \rv{windows.e3.working.cpu.mean} working (one core
is 100). The values of each system are read within that system (Section~\ref{sec:cost-method}); the table gives the
median, the 95th percentile and the maximum of the CPU, the final memory and its slope for each condition.

% Generated by manuscript/tools/results_tables.py from the result files; do not edit by hand.
\begin{sidewaystable}[p]
\caption{Runtime cost of the logger per system and condition: CPU of the logger (one core is 100) over the samples, its final resident memory (macOS) or working set (Windows) and the least-squares slope of it. The samplers of the two systems measure differently, so the values are read within each system and not compared between systems; --: not measured, or not to be quoted (section 6.10)}
\label{tab:cost}
\begin{tabular}{L{0.10}L{0.09}L{0.18}L{0.07}L{0.09}L{0.10}L{0.10}L{0.07}L{0.10}L{0.10}}
\toprule
System & Condition & Status & Mean & Median & 95th percentile & Maximum & Hours & Memory (MB) & Slope (MB/h) \\
\midrule
macOS & idle & measured & 0.32 & 0.32 & 0.35 & 0.38 & 7.99 & 5.0 & 0.00 \\
 & working & measured & 0.36 & 0.35 & 0.47 & 0.55 & 7.99 & 5.2 & -0.00 \\
Windows & idle & measured & 5.56 & 4.68 & 10.82 & 20.21 & 7.99 & 15.1 & 0.01 \\
 & working & measured & 5.37 & 4.66 & 10.84 & 15.51 & 7.99 & 16.3 & 0.09 \\
\bottomrule
\end{tabular}
\end{sidewaystable}

